\documentclass{article}
\usepackage[utf8]{inputenc}
\usepackage{amsmath}
\usepackage{graphicx}
\usepackage{listings}
\usepackage{hyperref}

\title{Práctica 1 Análisis Exploratorio de Datos de Nanopartículas con R}
\author{}
\date{}

\begin{document}





\begin{center}
\textbf{INGENIER\'IA EN NANOTECNOLOG\'IA \\
PROBABILIDAD Y ESTAD\'{I}STICA 2024-2025-I}
\vspace{.1cm}

Prof. Luis Jos\'{e} Yudico Anaya\\
 \hspace{.5cm}(19/09/2024, Práctica 1 )
\end{center}

\textbf{Análisis Exploratorio de Datos de Nanopartículas con R}

%\maketitle


\section*{Objetivos}
\begin{itemize}
    \item Familiarizar a los estudiantes con la carga y manipulaci\'on de grandes bases de datos en R.
    \item Aplicar conceptos de estad\'istica descriptiva para analizar las propiedades de las nanopart\'iculas.
    \item Visualizar los datos utilizando gr\'aficos relevantes.
    \item Interpretar los resultados obtenidos y extraer conclusiones preliminares.
\end{itemize}

\section*{Desarrollo de la Pr\'actica}

\subsection*{1. Carga de la Base de Datos}
\begin{itemize}
    \item Importar la base de datos a R utilizando funciones como \texttt{read.csv()}.
    \item Explorar las primeras filas y dimensiones de la base de datos con \texttt{head()}, \texttt{tail()}, \texttt{dim()}.
\end{itemize}

\subsection*{2. Estad\'istica Descriptiva}
\subsubsection*{Medidas de Tendencia Central}
\begin{itemize}
    \item Calcular la media, mediana y moda para variables num\'ericas como di\'ametro, concentraci\'on y potencial zeta. (recuerda el código de la moda cálculada en clase)
    \item Interpretar los resultados en el contexto de las nanopart\'iculas y el art\'iculo.
\end{itemize}

\subsubsection*{Medidas de Dispersi\'on}
\begin{itemize}
    \item Filtrando por Tipo de Nanopart\'icula en R.
    \item Calcular la varianza, desviaci\'on est\'andar, rango y coeficiente de variaci\'on para las variables num\'ericas de un tipo particular de nanopart\'icula utilizando funciones como:
    \begin{lstlisting}[language=R]
varianza <- var(data$variable)
desviacion_estandar <- sd(data$variable)
rango <- range(data$variable)
coef_variacion <- sd(data$variable) / mean(data$variable)
    \end{lstlisting}
    \item Explicar el significado del coeficiente de variaci\'on y su utilidad para comparar la dispersión entre variables.
\end{itemize}

\subsubsection*{Kurtosis}
\begin{itemize}
    \item Calcular la kurtosis para las variables num\'ericas utilizando la funci\'n \texttt{kurtosis()} desarrollada en clase.
    \item Interpretar si las distribuciones son leptoc\'urticas, mesoc\'urticas o platic\'urticas.
\end{itemize}

\subsection*{3. Visualizaci\'on de Datos}
\begin{itemize}
    \item \textbf{Histogramas:}
    \begin{itemize}
        \item Crear histogramas para variables num\'ericas como di\'ametro y concentraci\'on:
        \begin{lstlisting}[language=R]
 hist(data$diametro, main="Histograma de Di\'ametro",
 xlab="Di\'ametro", breaks=20)
        \end{lstlisting}
    \end{itemize}

    \item \textbf{Boxplots:}
    \begin{itemize}
        \item Crear boxplots para comparar la distribuci\'on de una variable num\'erica entre diferentes categor\'ias:
        \begin{lstlisting}[language=R]
boxplot(data$diametro ~ data$tipo_nanoparticula,
main="Boxplot de Di\'ametro por Tipo de Nanopart\'icula")
        \end{lstlisting}
    \end{itemize}

    \item \textbf{Gr\'aficos de Dispersi\'on:}
    \begin{itemize}
        \item Crear gr\'aficos de dispersi\'on para explorar la posible  relaci\'on entre dos variables num\'ericas:
        \begin{lstlisting}[language=R]
plot(data$diametro, data$concentracion, 
main="Di\'ametro vs Concentraci\'on",
xlab="Di\'ametro", ylab="Concentraci\'on")
        \end{lstlisting}
    \end{itemize}

    \item \textbf{Gr\'aficos de Barras:}
    \begin{itemize}
        \item Crear gr\'aficos de barras para variables categóricas:
        \begin{lstlisting}[language=R]
barplot(table(data$tipo_nanoparticula), 
main="Frecuencia de Tipos de Nanopart\'iculas", 
xlab="Tipo de Nanopart\'icula")
        \end{lstlisting}
    \end{itemize}
\end{itemize}

\subsection*{4. Interpretaci\'on de Resultados}
\begin{itemize}
    \item Interpretar los resultados obtenidos en los pasos anteriores.
    \item ¿Cu\'al es el tamaño promedio de las nanopart\'iculas?
    \item ¿Existe una gran variabilidad en la concentraci\'on?
    \item ¿Hay diferencias en el potencial zeta entre los diferentes tipos de nanopart\'iculas?
    \item ¿Cu\'al es la nanopart\'icula m\'as estudiada?
    \item ¿Cu\'ales son las nanopart\'iculas que se estudiaron en fecha m\'as reciente?
    \item ¿Cu\'ales son los tres primeros tipos de c\'elulas m\'as utilizados en an\'alisis de toxicidad en este art\'iculo?
    \item Agrega algo que haya llamado tu atención en tu an\'alisis
\end{itemize}

\section*{Consideraciones Adicionales}
\begin{itemize}
    \item Los datos de la pr\'actica son parte del art\'iculo \textit{Meta-Analysis of Nanoparticle Cytotoxicity via Data-Mining the Literature} y puede descargarse en \url{https://pubs.acs.org/doi/10.1021/acsnano.8b07562}. El archivo a utilizar es \texttt{nn8b07562\_si\_0012.csv}.
    \item \textbf{Existen varias formas de filtrar datos en R}, pero las m\'as comunes y eficientes son utilizando el paquete \texttt{dplyr} del \texttt{tidyverse}.
   % \item \textbf{Caso Insensible:} Si los nombres de los tipos de nanopart\'iculas no son sensibles a may\'usculas y min\'usculas, puedes utilizar funciones como \texttt{tolower()} o \texttt{toupper()} para normalizar los datos antes de filtrar.
    \item \textbf{Múltiples Condiciones:} Puedes combinar m\'ultiples condiciones utilizando operadores l\'ogicos como \texttt{\&} (y) y \texttt{|} (o). Por ejemplo:
    \begin{lstlisting}[language=R]
filter(Tipo_Nanoparticula = = "Organico" & Concentracion > 10)
    \end{lstlisting}
    \item \textbf{Otras Funciones \'Utiles:}
    \begin{itemize}
        \item \texttt{select()}: Selecciona columnas espec\'ificas.
        \item \texttt{arrange()}: Ordena los datos.
        \item \texttt{mutate()}: Crea nuevas variables.
        \item \texttt{group\_by()}: Agrupa los datos.
        \item \texttt{summarize()}: Calcula estad\'isticas de resumen por grupo.
    \end{itemize}
\end{itemize}

\end{document}

